\documentclass[aspectratio=169,usenames,dvipsnames]{beamer}
\beamertemplatenavigationsymbolsempty
% \documentclass{beamer}

% UNCOMMENT FOR HANDOUTS
% \usepackage{pgfpages}
% \pgfpagesuselayout{2 on 1}[landscape,letterpaper,border shrink=5mm]
% \setbeameroption{show notes}

\mode<presentation>
\usetheme{Boadilla}
\usecolortheme{spruce}
\usepackage{../LizStyle,ulem}
\usepackage{color}
\usepackage[color,matrix,arrow]{xy}
% \input xy
% \xyoption{all}
\usepackage{tikz}
\usepackage{tikz-cd}
\tikzset{commutative diagrams/diagrams={ampersand replacement=\&}}

\AtBeginSection{\frame{\sectionpage}}


% ==========Command for putting comments and removing them=======
% Visible:
\newcommand{\InstructorNotes}[1]{\textit{\color{orange}#1}}
\newcommand{\InstructorList}[1]{
\begin{itemize}
	\color{orange}
	#1
\end{itemize}
}
% Removed:
% \renewcommand{\InstructorNotes}[1]{}
% \renewcommand{\InstructorList}[1]{}


\newcommand{\bvfill}{\vskip0pt plus 1filll}

\usepackage{multimedia}
\usepackage{graphbox}

% \usepackage{algpseudocode}


\graphicspath{ {Figures/} {../Figures/}{../../Figures/} }



\begin{document}
%\pdfoptionpdfminorversion 5
%\logo{\includegraphics[height=1.5cm]{dukeshield.pdf}}
\title[]{Ch 7.4: Cubic splines}
\subtitle{Lecture 23 - CMSE 381}
\date{Fri, Nov 1, 2024}



\institute[MSU-CMSE]{Michigan State University\\
			::\\
          Dept of Computational Mathematics, Science \& Engineering}
\author[Dr.~Munch]{Prof.~Elizabeth Munch}

\begin{frame}
 \titlepage


\end{frame}
%--------------------------
\begin{frame}[c]{Announcements}

\begin{columns}
\column{.45\textwidth}
\centering

	\textbf{Last time:}
	\begin{itemize}
		\item 7.2 Step functions 
		\item 7.3 Basis functions
	\end{itemize}

	\vspace{.2in}

	\textbf{This lecture:}
	\begin{itemize}
	\item 7.4 Cubic splines
\end{itemize}
\InstructorList{
  \item For next time, fix this to deal with the fact that the book basis functions for cubic splines is not the same as the basis used by \texttt{scikitlearn}
}


\textbf{Announcements:}
	\begin{itemize}
		\item Homework \# 6 is  due Sunday 
		\item Projects

	\end{itemize}


\column{.4\textwidth}
\centering

\includegraphics[width=\textwidth]{Day01-Schedule-Test3.png}
\end{columns}

\end{frame}
%--- Next Frame ---%



%=================
\section{Last time}
%=================



\begin{frame}[c]{Polynomial regression}

\begin{columns}
\column{.45\textwidth}
\centering
\begin{equation*}
	y_i = \beta_0 + \beta_1x_1 + \beta_2x_i^2 + \cdots + \beta_dx_i^d + \e_i
\end{equation*}

\column{.4\textwidth}
\centering
\includegraphics[width = \textwidth, align = c]{Fig7-1a}

\end{columns}

\end{frame}
%--- Next Frame ---%
\begin{frame}[t]{Step function regression }
\begin{columns}
\column{.5\textwidth}
\centering
\includegraphics[width = .8\textwidth, align = c]{Eq7-4}
\begin{equation*}
	y_i = \beta_0 + \beta_1C_1(x_i) + \beta_2 C_2(x_i) + \cdots + \beta_K C_K(x_i) + \e_i
\end{equation*}
\column{.4\textwidth}
\centering
\includegraphics[width = \textwidth, align = c]{Fig7-2a}
\end{columns}
\end{frame}
%--- Next Frame ---%

\begin{frame}[c]{Classification version}

\begin{equation*}
	\Pr(y_i>250 \mid x_i) =
	% \frac{\exp(\beta_0 + \beta_1x_i + \cdots + \beta_dx_i^d)}
	    %  {1+\exp(\beta_0 + \beta_1x_i + \cdots + \beta_dx_i^d)}
\end{equation*}

\vspace{-.2in}
\begin{columns}
\column{.35\textwidth}
\centering
\footnotesize
\begin{equation*}
	% \Pr(y_i>250 \mid x_i) =
	\frac{\exp(\beta_0 + \beta_1x_i + \cdots + \beta_dx_i^d)}
	     {1+\exp(\beta_0 + \beta_1x_i + \cdots + \beta_dx_i^d)}
\end{equation*}
\includegraphics[width = .8\textwidth, align = c]{Fig7-1b}
\column{.5\textwidth}
\centering

\footnotesize
\begin{equation*}
	% \Pr(y_i>250 \mid x_i) =
	\frac{\exp(\beta_0 + \beta_1C_1(x_i) + \beta_2 C_2(x_i) + \cdots + \beta_K C_K(x_i))}
	     {1+\exp(\beta_0 + \beta_1C_1(x_i) + \beta_2 C_2(x_i) + \cdots + \beta_K C_K(x_i))}
\end{equation*}
\includegraphics[width = .55\textwidth, align = c]{Fig7-2b}
\end{columns}

\end{frame}
%--- Next Frame ---%


\begin{frame}[t]{Basis Functions Setup}
Polynomial and piecewise-constant regression models are special cases of a \textit{basis function} approach.

\begin{equation*}
	y_i = \beta_0 + \beta_1b_1(x_i) + \beta_2 b_2(x_i) + \cdots + \beta_K b_K(x_i) + \e_i
\end{equation*}

\begin{columns}
\column{.45\textwidth}
\centering
\InstructorList{
  \item Pick a collection of basis functions $b_1(X),\cdots,b_K(X)$
	\item Use least squares to figure out the constants

}

\column{.45\textwidth}
\centering

\InstructorList{
	\item Lots of possible options for these
	\InstructorList{
	  \item Some examples are wavelets or Fourier series
		\item Next section is we'll talk about regression splines
	}
}


\end{columns}

\end{frame}
%--- Next Frame ---%
%==========
\section{Regression Splines}
%==========


\begin{frame}[c]{Piecewise polynomials}

\begin{columns}
\column{.45\textwidth}
\centering
\begin{itemize}
	\item Fit a polynomial regression
\begin{equation*}
	y_i = \beta_0 + \beta_1x_1 + \beta_2x_i^2 + \cdots + \beta_dx_i^d + \e_i
\end{equation*}
\item Let the $\beta_i$'s be different at different locations of the range.
\end{itemize}
\column{.45\textwidth}
\centering
\InstructorList{
  \item Points where the coefficients change are called "knots"
	\item A piecewise cubic with no knots is just standard cubic polynomial
	\item More knots gives more flexible polynomial
	\item Piecewise constant of previous section is degree 0 version of this
}

\end{columns}

\end{frame}
%--- Next Frame ---%

\begin{frame}[t]{Building up to cubic splines}
	\begin{columns}
	\column{.3\textwidth}
	\centering
	
\includegraphics[width = \textwidth, align = c]{Fig7-3a}
\InstructorNotes{Piecewise cubic but discontinuous}
	
	\column{.3\textwidth}
	\centering
\includegraphics[width = \textwidth, align = c]{Fig7-3b}
\InstructorNotes{Continuous but pointy join}

	\column{.3\textwidth}
	\centering

\includegraphics[width = \textwidth, align = c]{Fig7-3c}
\InstructorNotes{Smooth join. Same value plus same first and second derivatives.}
	\end{columns}
\bvfill 

	\centering
\includegraphics[width = .5\textwidth, align = c]{Eq7-8b.png}
\end{frame}
%--- Next Frame ---%

%--- Next Frame ---%
\begin{frame}[t]{Cubic splines: degrees of freedom}

\begin{columns}
\column{.45\textwidth}
\centering
\begin{equation*}
	f(x) = 
	\begin{cases}
		\beta_0^1 + \beta_1^1 x + \beta_2^1 x^2 + \beta_3^1 x^3 & x<c \\
		\beta_0^2 + \beta_1^2 x + \beta_2^2 x^2 + \beta_3^2 x^3 & x>c
	\end{cases}
\end{equation*}

% \begin{itemize}
% 	\item Fit piecewise polynomial
% 	\item Require continuity at knots
% 	\item Require the first and second derivatives to be continuous at knots
% \end{itemize}
\InstructorList{
	\item 8 original degrees of freedom, plus three constraints (continuous, 1st deriv continuous, 2nd deriv continuous) leaves 5 degrees of freedom
	\item cubic spline with $K$ knots uses $4+K$ degrees of freedom
}
\column{.45\textwidth}
\centering
\includegraphics[width = \textwidth, align = c]{Fig7-3c}

\end{columns}

\end{frame}
%--- Next Frame ---%

% \begin{frame}[c]{Generalized splines}

% \begin{columns}
% \column{.45\textwidth}
% \centering
% \InstructorList{
%   \item Can do this for any degree $d$ polynomial
% 	\item Require continuity up to degree $d-1$ at each knot
% 	\item Linear then is just requiring continuity
% }

% \column{.45\textwidth}
% \centering
% \includegraphics[width = \textwidth, align = c]{Fig7-3d}
% \end{columns}

% \end{frame}
% %--- Next Frame ---%

\begin{frame}[c]{Spline basis representation}

\begin{columns}
\column{.65\textwidth}
\centering
Want to pick $b_i$ so that we represent a cubic spline with $K$ knots as
\begin{equation*}
	y_i = \beta_0 + \beta_1b_1(x_i) + \beta_2 b_2(x_i) + \cdots + \beta_{K+3} b_{K+3}(x_i) + \e_i
\end{equation*}
\end{columns}

\InstructorNotes{Warning: There is no unique way to do this. The book has a different way. We're going to use scikit learn's way.}

\end{frame}
%--- Next Frame ---%


\begin{frame}[t]{Version 1: Truncated power basis function}

\begin{columns}
\column{.45\textwidth}
\centering

\begin{equation*}
h(x,z) = (x-z)^3_+ =
\begin{cases}
	(x-z)^3 & \text{if } x>z \\
	0 & else
\end{cases}
\end{equation*}
Desmos link:
\url{https://www.desmos.com/calculator/esucuulbgj}
\column{.45\textwidth}
\centering
\InstructorList{
  \item Book uses $z=\xi$ notation, but i refuse
	\item Draw $h(x,2)$
}

\end{columns}

\end{frame}
%--- Next Frame ---%
\begin{frame}[t]{The (book) basis for cubic splines }

\begin{columns}
\column{.45\textwidth}
\centering
Given knots at $z_1,\cdots,z_K$
\begin{itemize}
	\item $X$
	\item $X^2$
	\item $X^3$
	\item $h(X,z_1)$
	\item $h(X,z_2)$
	\item $\vdots$
	\item $h(X,z_K)$
\end{itemize}
\column{.45\textwidth}
\centering


\InstructorList{
  \item Adding $h(X,z_i)$ leads to discontinuity at most at the third derivative at $z_i$
	\item So use the $3+K$ functions listed here, result is $K+4$ coefficients once you get the intercept.
	\item Matches with our $K+4$ degrees of freedom from earlier
}

\end{columns}

\bvfill

\begin{equation*}
	f(X) = 	
	\beta_0 + 
	\beta_1 X + 
	\beta_2 X^2 + 
	\beta_3 X^3 + 
	\beta_4 h(X,z_1) + 
	\beta_5 h(X,z_2) + 
	\cdots + 
	\beta_{k+3}  h(X,z_K)
\end{equation*}

\end{frame}
%--- Next Frame ---%
\begin{frame}[t]{Version 2: B-spline basis function}
	\subtitle{sklearn's basis}

	\centering
	\includegraphics[width = .8\textwidth]{B-Spline_Basis_sklearn.png}

	\InstructorNotes{Similar idea. Take some functions that are all continuous and differentiable everywhere and add them up so you end up with a function that is continuous and differentiable everywhere.}
\end{frame}
%--- Next Frame ---%

{
\setbeamercolor{background canvas}{bg=Apricot!25}
\begin{frame}[t]{Coding example}
\begin{columns}
\column{.45\textwidth}
\centering
\column{.45\textwidth}
\centering
\end{columns}
\end{frame}
}
%--- Next Frame ---%

\begin{frame}[t]{Notes on cubic splines }

\begin{columns}
\column{.45\textwidth}
\centering

\InstructorList{
  \item Human eyes can't detect the discontinuities so often cubic splines are a good way to go
	\item Still have problem of high variance at the outer range of predictors (big or small $X$)
	\item \textit{Natural spline:} regression spline plus boundary constraints where function is required to be linear at the boundary
}

\column{.45\textwidth}
\centering

\includegraphics[width = \textwidth]{Fig7-4}
\end{columns}

\end{frame}
%--- Next Frame ---%
\begin{frame}[t]{Where to put the knots?}

\begin{columns}

\column{.65\textwidth}
\centering

\includegraphics[width = \textwidth]{Fig7-5}
\column{.25\textwidth}
\centering

\InstructorList{
  \item Plan A: Uniform placement
	\item Plan B: Data driven. Example here puts it at quantiles of the data
}
\end{columns}

\end{frame}
%--- Next Frame ---%
\begin{frame}[t]{How many knots to use?}
{When in doubt, Cross-Validate.}
\begin{columns}
\column{.45\textwidth}
\centering

\InstructorList{
  \item Do CV ($k$-fold probably) on data for each number of knots
	\item Figure at right is the MSE resulting from 10-fold CV
	\item Note that clearly linear is not enough, after that 3-4 degrees of freedom for the cubic spline are enough
}

\column{.45\textwidth}
\includegraphics[width = \textwidth]{Fig7-6b}
\centering

\end{columns}
\end{frame}
%--- Next Frame ---%
\begin{frame}[t]{Cubic splines vs Polynomial Regression}

\begin{columns}
\column{.45\textwidth}
\centering
\includegraphics[width = \textwidth]{Fig7-7}
\column{.45\textwidth}
\centering

\InstructorList{
  \item Left: natural cubic spline with 15 degrees of freedom vs degree 15 polynomial
	\item Polynomial flexibility means ugly boundary behavior
	\item Splines let you add flexibility as more knots rather than more degrees, so allows for flexibility where the function is changing more rapidly
}

\end{columns}

\end{frame}
%--- Next Frame ---%


\begin{frame}[t]{Next time}

\includegraphics[align = c, width = .35\textwidth]{Day01-Schedule-Test3}


\end{frame}
%--- Next Frame ---%


\end{document}
